\documentclass[12pt,a4paper]{article}

\usepackage[spanish,es-tabla]{babel}
\usepackage[T1]{fontenc}
\usepackage[utf8]{inputenc}
\usepackage{lmodern}
\usepackage{geometry}
\usepackage{setspace}
\usepackage{array}
\usepackage{longtable}
\usepackage{booktabs}
\usepackage{hyperref}
\usepackage{enumitem}
\usepackage{tikz}
\usepackage{float}
\usepackage{graphicx}
\usepackage{xcolor}
\usepackage{listings}

\usetikzlibrary{arrows.meta,positioning,shapes.geometric,fit}

\definecolor{ullblue}{RGB}{0,51,102}
\definecolor{codegray}{RGB}{245,245,245}
\definecolor{codeblue}{RGB}{0,51,102}
\definecolor{codegreen}{RGB}{0,110,55}

\lstdefinestyle{pythoncode}{
  language=Python,
  backgroundcolor=\color{codegray},
  basicstyle=\ttfamily\small,
  keywordstyle=\color{codeblue}\bfseries,
  commentstyle=\color{codegreen},
  stringstyle=\color{orange!70!black},
  numbers=left,
  numberstyle=\tiny\color{gray},
  stepnumber=1,
  numbersep=8pt,
  frame=single,
  rulecolor=\color{gray!45},
  breaklines=true,
  showstringspaces=false,
  tabsize=4,
  keepspaces=true,
  captionpos=b
}

\geometry{margin=2.5cm}
\setstretch{1.15}
\setlist[itemize]{leftmargin=1.3cm}
\setlist[enumerate]{leftmargin=1.1cm}
\setlength{\tabcolsep}{4pt}
\renewcommand{\arraystretch}{1.12}
\newcolumntype{P}[1]{>{\raggedright\arraybackslash}p{#1}}

\begin{document}

\begin{titlepage}
\centering

% ── Franja superior ────────────────────────────────────────────
{\color{ullblue}\rule{\linewidth}{1.5pt}}\\[0.4cm]
{\large\scshape Universidad de La Laguna}\\[0.15cm]
{\normalsize Escuela Superior de Ingenier\'ia y Tecnolog\'ia}\\[0.15cm]
{\normalsize Grado en Ingenier\'ia Electr\'onica Industrial y Autom\'atica}\\[0.3cm]
{\color{ullblue}\rule{\linewidth}{0.6pt}}

\vspace{2cm}

% ── Asignatura ─────────────────────────────────────────────────
{\large Sistemas Ciberf\'isicos\quad|\quad Curso 2025/26}

\vspace{2cm}

% ── Título principal ───────────────────────────────────────────
{\color{ullblue}\rule{0.6\linewidth}{0.4pt}}\\[0.6cm]
{\LARGE\bfseries Documento T\'ecnico Final}\\[0.5cm]
{\Large Sistema ciberf\'isico:\\[0.2cm] 3 en raya con Ned-2}\\[0.6cm]
{\color{ullblue}\rule{0.6\linewidth}{0.4pt}}

\vfill

% ── Autor e info ───────────────────────────────────────────────
\begin{tabular}{rl}
  \textbf{Autor:}  & Pablo Navarro Dom\'inguez \\
  \textbf{Fecha:}  & 14 de mayo de 2026 \\
\end{tabular}

\vspace{1.2cm}
{\color{ullblue}\rule{\linewidth}{1.5pt}}

\end{titlepage}

\tableofcontents
\newpage

\section{Resumen}
Este documento t\'ecnico final consolida el desarrollo del sistema ciberf\'isico de 3 en raya con Ned-2, integrando los avances recogidos en los dos documentos t\'ecnicos previos y aline\'andolos con el enunciado del proyecto. El objetivo funcional del sistema es ejecutar una partida f\'isica completa frente a un rival humano, combinando supervisi\'on en tiempo real, manipulaci\'on segura, detecci\'on visual del tablero y comunicaci\'on remota entre nodos de control.

La soluci\'on se apoya en componentes concretos implementados y validados: un n\'ucleo Python para control de partida, movimiento y visi\'on en el robot, junto con bloques MATLAB/Simulink en el PC central para emisi\'on y recepci\'on de mensajes TCP/IP.
\par
La implementaci\'on se estructura en dos capas de decisi\'on y ejecuci\'on: el PC central, responsable del control supervisado mediante m\'aquina de estados en Simulink, y el PC del robot, donde se ejecuta la l\'ogica Python de movimiento, visi\'on, inteligencia artificial y mensajer\'ia. El robot mantiene adem\'as una segunda conexi\'on TCP/IP opcional hacia otro robot Ned-2 rival, lo que permite extender el escenario de juego a una partida m\'aquina contra m\'aquina en la que cada robot recibe los movimientos del contrario. Esta separaci\'on ha permitido mantener trazabilidad de eventos, reducir acoplamiento entre subsistemas y facilitar la verificaci\'on incremental por hitos.
\par
En la versi\'on final del proyecto el sistema dispone de funcionamiento operativo completo en ambos modos de juego: el modo manual, en el que el PC central emite los movimientos a ejecutar, y el modo autom\'atico, en el que el propio robot calcula su jugada mediante una heur\'istica que prioriza la victoria inmediata y, en su defecto, la ocupaci\'on del centro o una posici\'on aleatoria libre. El protocolo de mensajes TCP/IP cubre el ciclo completo de partida, incluyendo la emisi\'on sistem\'atica de los mensajes de cierre \texttt{FIN VIC}, \texttt{FIN DER} y \texttt{FIN EMP} en todos los caminos de terminaci\'on posibles. La detecci\'on visual del tablero, las secuencias de movimiento pick-and-place y el flujo de sincronizaci\'on de turnos mediante la entrada digital \texttt{DI1} se encuentran igualmente operativos.

\section{Alcance del documento final}
El alcance de este informe incluye la descripci\'on t\'ecnica de la soluci\'on desarrollada en el PC central, el PC de control del robot, el subsistema de visi\'on y la capa de comunicaciones. Tambi\'en cubre las decisiones de dise\~no adoptadas para satisfacer las restricciones del enunciado, en particular las asociadas a seguridad cinem\'atica, sincronizaci\'on de turnos y capacidad de supervisi\'on remota.

Adem\'as, se incorporan detalles de hardware y software que ya aparec\'ian en los dos documentos previos: la configuraci\'on de visi\'on del Ned-2, la posici\'on home como estado seguro de inicio, las limitaciones temporales de recogida y transporte, el uso de mensajes de texto terminados en salto de l\'inea y el papel de la IA en el modo autom\'atico.
\par
Desde el punto de vista de verificaci\'on, el documento organiza el estado de cumplimiento de requisitos, indicando qu\'e aspectos ya fueron validados en entregables anteriores y qu\'e aspectos se validan en este documento. Se mantiene, adem\'as, una estructura de evidencias compatible con la trazabilidad exigida por el cliente.

\begin{table}[H]
\centering
\footnotesize
\begin{tabular}{P{0.18\textwidth} P{0.26\textwidth} P{0.26\textwidth} P{0.20\textwidth}}
\toprule
\textbf{Hito} & \textbf{Objetivo de proyecto} & \textbf{Cobertura documental actual} & \textbf{Estado} \\
\midrule
H1 & Aprobaci\'on del documento de requisitos & Referencia de requisitos y matriz de validaci\'on & Completado \\
H2 & Verificaci\'on del PC central & Arquitectura de control y protocolo Simulink & Completado \\
H3 & Verificaci\'on del robot manipulador & Control de movimiento, visi\'on y seguridad & Completado \\
H4 & Validaci\'on del sistema ciberf\'isico & Procedimientos y evidencias de cierre & Completado \\
\bottomrule
\end{tabular}
\caption{Cobertura del documento respecto a los hitos del proyecto.}
\end{table}

\section{Descripci\'on t\'ecnica de la soluci\'on final}

\subsection{Arquitectura general del sistema}
La arquitectura final implementa un esquema distribuido de supervisi\'on y ejecuci\'on. El PC central act\'ua como nodo de control de alto nivel y concentra la interfaz de usuario, la m\'aquina de estados y la coordinaci\'on de la partida. El PC del robot ejecuta la capa de campo: control de movimiento, percepci\'on visual, c\'alculo de jugada en modo autom\'atico y notificaci\'on de eventos al sistema central.
\par
El ciclo operativo comienza con la selecci\'on de modo y la verificaci\'on de precondiciones de seguridad (home + Start), contin\'ua con la alternancia de turnos y finaliza cuando se detecta una condici\'on de victoria, derrota o empate. Esta estructura se dise\~n\'o para satisfacer las exigencias del enunciado sobre supervisi\'on continua del estado y control remoto del robot.

\begin{figure}[h]
\centering
\begin{tikzpicture}[
  box/.style={draw, rounded corners, align=center, minimum width=3.4cm, minimum height=1.1cm, fill=gray!8, font=\small},
  arrow/.style={-Latex, thick},
  lab/.style={font=\small, fill=white, inner sep=1.5pt}
]
\node[box] (pc) at (0,0) {PC central\\Interfaz + Simulink};
\node[box] (pcr) at (8.4,0) {PC robot\\Python de control};
\node[box] (vis) at (0,-4.0) {Visi\'on artificial\\Detecci\'on tablero};
\node[box] (ned) at (8.4,-4.0) {Robot Ned-2\\Movimiento + audio};

% PC central -> PC robot
\draw[arrow] (pc.east) -- node[lab,above,pos=0.5]{INI, MOV} (pcr.west);

% PC robot -> PC central (ruta superior para evitar solape)
\draw[arrow] (pcr.north) |- (8.4,1.8) -- (0,1.8) -| node[lab,above,pos=0.5]{DEC, ACT, FTJ, FTU, FIN} (pc.north);

% PC robot -> Robot
\draw[arrow] (pcr.south) -- node[lab,right,pos=0.5,xshift=2pt]{Comandos API} (ned.north);

% Robot -> Vision
\draw[arrow] (ned.west) -- node[lab,below,pos=0.5]{Imagen} (vis.east);

% Vision -> PC robot (ruta lateral para no cruzar etiquetas)
\draw[arrow] (vis.north) |- (0,-2.55) -- (8.4,-2.55) -| (pcr.south);
\node[lab] at (4.2,-3.65) {Estado visual};
\end{tikzpicture}
\caption{Gr\'afica de arquitectura funcional y flujo de informaci\'on del sistema.}
\end{figure}

\begin{table}[H]
\centering
\footnotesize
\begin{tabular}{P{0.22\textwidth} P{0.25\textwidth} P{0.45\textwidth}}
\toprule
\textbf{Subsistema} & \textbf{Tecnolog\'ia principal} & \textbf{Funci\'on principal en el sistema} \\
\midrule
PC central & Matlab/Simulink & Supervisar estados, coordinar turnos y presentar estado de partida al usuario. \\
PC robot & Python + sockets & Ejecutar l\'ogica de juego, control de movimiento y comunicaci\'on con el PC central. \\
Robot Ned-2 & API pyniryo & Manipular piezas en tablero siguiendo restricciones cinem\'aticas de seguridad. \\
Visi\'on artificial & OpenCV + c\'amara del robot & Detectar estado del tablero rival tras su turno y actualizar partida. \\
\bottomrule
\end{tabular}
\caption{Resumen de responsabilidades t\'ecnicas por subsistema.}
\end{table}

\subsection{Subsistemas software desarrollados}
Los ficheros desarrollados para el PC del robot y el PC central permiten distinguir claramente las capas de responsabilidad del sistema. En Python, ComRobot.py concentra el bucle de partida, control de turnos, sockets y secuencias de movimiento; Ned2\_ULL.py agrupa utilidades de configuraci\'on y persistencia de poses; y VisionRobot.py implementa la detecci\'on del estado del tablero. En Simulink/MATLAB, el intercambio de datos y su adaptaci\'on al modelo se materializan mediante bloques de emisi\'on, recepci\'on e interpretaci\'on de se\~nales.

\begin{table}[H]
\centering
\footnotesize
\begin{tabular}{P{0.22\textwidth} P{0.20\textwidth} P{0.48\textwidth}}
\toprule
\textbf{Archivo} & \textbf{Tipo} & \textbf{Responsabilidad principal} \\
\midrule
ComRobot.py & Python & Bucle principal de partida, secuenciaci\'on de movimientos, gesti\'on de turnos y mensajes de estado. \\
Ned2\_ULL.py & Python & Utilidades de persistencia de poses, restauraci\'on de configuraci\'on y apoyo al robot. \\
VisionRobot.py & Python & Segmentaci\'on visual del tablero y extracci\'on de casillas ocupadas para actualizar el estado de juego. \\
SimulinkEmisor.m & MATLAB/Simulink & Formateo de comandos INI/MOV en tramas \texttt{uint8[1x32]} para el bloque TCP/IP Send. \\
SimulinkReceptor.m & MATLAB/Simulink & Recepci\'on con buffer persistente, parseo de DEC/ACT/FTU/FTJ/FIN y generaci\'on de pulsos y estado de tablero. \\
InterpretarTablero.m & MATLAB/Simulink & Descomposici\'on del vector de tablero 1x9 en nueve se\~nales escalares para l\'ogica de estados/visualizaci\'on. \\
MensajeTurno.m & MATLAB/Simulink & Generaci\'on de mensaje ASCII de turno (\texttt{Espera} / \texttt{Te toca jugar}) para interfaz/supervisi\'on. \\
\bottomrule
\end{tabular}
\caption{Relaci\'on de componentes software desarrollados y su funci\'on t\'ecnica.}
\end{table}

\subsection{Caracter\'isticas principales y decisiones de dise\~no por subsistema}

\subsubsection{PC central (Simulink + control de partida)}
El PC central implementa el control supervisado en tiempo real solicitado en el enunciado. La m\'aquina de estados en Simulink secuencia las fases de espera, orden de movimiento, confirmaci\'on de ejecuci\'on y evaluaci\'on de continuidad de partida. Esta representaci\'on permite validar transiciones de forma expl\'icita y facilita la interpretaci\'on del estado por parte del operador.
\par
La interfaz de comunicaci\'on se realiza mediante funciones MATLAB que convierten comandos y eventos entre se\~nales de simulaci\'on y tramas de protocolo. Se adopt\'o una sem\'antica de pulsos de un ciclo para eventos discretos (DEC, FTJ, FTU, FIN), evitando disparos m\'ultiples y mejorando el acoplamiento temporal con el modelo de estados.

Los bloques MATLAB incluidos en el repositorio cubren de forma expl\'icita la frontera de comunicaci\'on con el robot. SimulinkEmisor.m genera siempre un vector fijo de 32 bytes de tipo \texttt{uint8}, lo que simplifica su conexi\'on con el bloque TCP/IP Send de Simulink y evita variabilidad de tama\~no en tiempo de ejecuci\'on. El bloque soporta mensajes \texttt{INI MAN}, \texttt{INI AUT}, \texttt{INI COM} y \texttt{MOV origen destino}, incluyendo la traducci\'on de identificadores de almac\'en 11--13 a las cadenas \texttt{alm1}--\texttt{alm3}. Tambi\'en contempla un modo de emisi\'on de \texttt{MOV} sin argumentos para escenarios de sincronizaci\'on o prueba.

Por su parte, SimulinkReceptor.m implementa un buffer persistente de recepci\'on, acumula bytes procedentes del bloque TCP/IP Receive y procesa l\'ineas completas terminadas en CR/LF. El dise\~no desacopla recepci\'on y decodificaci\'on, permitiendo absorber fragmentaci\'on de red sin perder mensajes. Sobre ese buffer se reconocen los comandos \texttt{DEC}, \texttt{ACT}, \texttt{FTU}, \texttt{FTJ} y \texttt{FIN}. Para los eventos discretos se utilizan contadores persistentes que mantienen pulsos de 100 ms equivalentes en muestras, calculados a partir del periodo \texttt{Ts} del modelo. De esta forma, el resto de la m\'aquina de estados recibe se\~nales limpias y temporizadas.

En el caso del mensaje \texttt{ACT}, el bloque receptor transforma la cadena de nueve celdas del tablero en un vector num\'erico de longitud 9 apto para tratamiento en Simulink. Esta conversi\'on permite que la l\'ogica del PC central trabaje directamente con un estado discreto del tablero, sin necesidad de parseado textual adicional. El mensaje \texttt{FIN} actualiza adem\'as la salida \texttt{RES} con etiquetas de texto (\texttt{victoria}, \texttt{derrota} o \texttt{empate}), lo que simplifica la integraci\'on con estados finales y paneles de supervisi\'on del modelo.

\begin{figure}[h]
\centering
\begin{tikzpicture}[
  stepbox/.style={draw, rounded corners, minimum width=2.8cm, minimum height=0.9cm, align=center, fill=gray!10, font=\small},
    link/.style={-Latex, thick}
]
\node[stepbox] (h) {Estado de reposo\\robot en home};
\node[stepbox, below=0.55cm of h] (st) {Start + selecci\'on de modo};
\node[stepbox, below=0.55cm of st] (tr) {Turno del robot\\ejecuci\'on de jugada};
\node[stepbox, below=0.55cm of tr] (tu) {Turno del rival\\actualizaci\'on por visi\'on};
\node[stepbox, below=0.55cm of tu] (ev) {Evaluaci\'on\\FIN / continuaci\'on};

\draw[link] (h) -- (st);
\draw[link] (st) -- node[right,pos=0.5]{INI} (tr);
\draw[link] (tr) -- node[right,pos=0.5]{FTJ} (tu);
\draw[link] (tu) -- node[right,pos=0.5]{ACT / FTU} (ev);
\draw[link] (ev.east) -- ++(1.0,0) |- node[pos=0.25,right]{si no FIN} (tr.east);
\end{tikzpicture}
\caption{Gr\'afica de ciclo de turno en la m\'aquina de estados del sistema.}
\end{figure}

\begin{table}[H]
\centering
\footnotesize
\begin{tabular}{P{0.24\textwidth} P{0.18\textwidth} P{0.48\textwidth}}
\toprule
\textbf{Bloque Simulink} & \textbf{Interfaz} & \textbf{Comportamiento implementado} \\
\midrule
SimulinkEmisor.m & Salida \texttt{uint8[1x32]} & Genera tramas de longitud fija para \texttt{INI} y \texttt{MOV}, con terminador de l\'inea y compatibilidad directa con TCP/IP Send. \\
SimulinkReceptor.m & Entradas \texttt{data,count,Ts} & Acumula bytes en buffer persistente, procesa mensajes por l\'inea y expone pulsos discretos y estado de tablero. \\
SimulinkReceptor.m & Salidas \texttt{DEC, FTU, FTJ, FIN, RES, tablero} & Entrega eventos discretos temporizados y un vector de 9 posiciones para la l\'ogica de control central. \\
\bottomrule
\end{tabular}
\caption{Resumen funcional del c\'odigo MATLAB/Simulink incluido en el PC central.}
\end{table}

\subsubsection{Subsistema robot (control de movimiento)}
El control de movimiento del Ned-2 est\'a implementado en ComRobot.py, donde se integra la recepci\'on de comandos, la planificaci\'on de secuencias pick-and-place y la notificaci\'on de estado. El sistema soporta modos manual y autom\'atico sobre una arquitectura com\'un de turnos, lo que permite conservar el mismo ciclo de ejecuci\'on independientemente de qui\'en decide la jugada. Adem\'as de la conexi\'on principal con el PC central, el c\'odigo contempla una segunda conexi\'on TCP/IP opcional hacia otro robot Ned-2, lo que habilita un escenario de juego m\'aquina contra m\'aquina en el que cada robot act\'ua como rival del otro.
\par
Se ha adoptado una estrategia de seguridad basada en desacoplo cinem\'atico Z/X-Y, tal como exige el enunciado. El robot primero se posiciona en una pose de aproximaci\'on, desciende verticalmente para recoger/dejar, y retorna al plano de aproximaci\'on antes de cualquier desplazamiento horizontal. Esta l\'ogica se complementa con limitaci\'on de velocidad (valor 50) y retorno a la pose de visi\'on general para el cierre de ciclo. El fin de turno rival se dispara mediante entrada digital DI1, tratada como evento de pulsador f\'isico.

\begin{table}[H]
\centering
\footnotesize
\begin{tabular}{P{0.28\textwidth} P{0.18\textwidth} P{0.42\textwidth}}
\toprule
\textbf{Elemento del Ned-2} & \textbf{Dato t\'ecnico} & \textbf{Uso en el proyecto} \\
\midrule
Grados de libertad & 6 & Permiten posicionar el efector final sobre el tablero y en la zona de reserva. \\
Sistema de visi\'on & C\'amara RGB integrada & Captura el tablero para detectar movimientos del rival. \\
Efector final & Pinza de dedos & Recogida y colocaci\'on de fichas. \\
Conectividad del robot & TCP/IP / rosbridge & Acceso mediante API pyniryo desde Python. \\
L\'imite de velocidad & 50 en la configuraci\'on usada & Garantizar movimientos suaves y seguros. \\
\bottomrule
\end{tabular}
\caption{Resumen del hardware y configuraci\'on funcional del Ned-2 empleada en el proyecto.}
\end{table}

\subsubsection{Sistema de visi\'on}
El subsistema de visi\'on utiliza la c\'amara integrada en el robot para actualizar el estado del tablero tras el turno rival. El pipeline de procesamiento incluye normalizaci\'on geom\'etrica del workspace, segmentaci\'on por color en HSV y asignaci\'on de detecciones a la rejilla 3x3. El resultado se publica como posiciones discretas en el protocolo de comunicaci\'on.
\par
Una decisi\'on de dise\~no especialmente relevante ha sido separar fuentes de verdad: el estado de fichas propias se mantiene por l\'ogica interna, mientras que el estado rival se obtiene por visi\'on. Esta separaci\'on reduce sensibilidad a errores de detecci\'on de las propias piezas y mejora la consistencia del ciclo de decisi\'on autom\'atica.

Durante el desarrollo se incorpor\'o una herramienta interactiva de calibraci\'on HSV mediante trackbars para facilitar ajustes seg\'un condiciones de iluminaci\'on. Esta funcionalidad fue especialmente importante en las pruebas iniciales, ya que el documento de requisitos no solo exige detectar la partida sino hacerlo de forma fiable y repetible en campo.

\subsubsection{Comunicaci\'on TCP/IP}
La comunicaci\'on remota se ha implementado sobre TCP/IP con mensajes ASCII terminados en fin de l\'inea. Esta elecci\'on prioriza simplicidad de depuraci\'on y compatibilidad directa entre Python y Simulink, manteniendo suficiente expresividad para incluir acci\'on solicitada, estado de proceso y estado de partida, tal como especifica el enunciado.
\par
En la direcci\'on PC central $\rightarrow$ robot se transmiten \'ordenes de inicializaci\'on y movimiento; en la direcci\'on opuesta se reportan eventos de decisi\'on, actualizaci\'on visual y cierre de turno. El receptor Simulink ya contempla mensaje FIN con resultado, pero su emisi\'on en el bucle principal Python se considera tarea abierta de cierre funcional.
En la direcci\'on PC central $\rightarrow$ robot se transmiten \'ordenes de inicializaci\'on y movimiento; en la direcci\'on opuesta se reportan eventos de decisi\'on, actualizaci\'on visual y cierre de turno. En la versi\'on actual de \texttt{ComRobot.py}, los mensajes de cierre \texttt{FIN VIC}, \texttt{FIN DER} y \texttt{FIN EMP} se emiten de forma sistem\'atica en todos los caminos de terminaci\'on de partida (victoria propia, derrota o empate), quedando alineados con el parser del receptor Simulink.

\begin{table}[H]
\centering
\footnotesize
\begin{tabular}{P{0.25\textwidth} P{0.26\textwidth} P{0.39\textwidth}}
\toprule
\textbf{Estructura del mensaje} & \textbf{Contenido} & \textbf{Interpretaci\'on} \\
\midrule
IDN arg1 arg2 ...\textbackslash n & Identificador + argumentos & Formato gen\'erico legible y compatible con parser por l\'inea. \\
INI modo\textbackslash n & Modo = MAN / AUT / COM & Inicio de partida en el modo seleccionado. \\
MOV origen destino\textbackslash n & Origen y destino & Orden de recogida y colocaci\'on en modo manual. \\
DEC\textbackslash n & Evento discreto & Confirmaci\'on de jugada decidida en modo autom\'atico. \\
ACT cadena9\textbackslash n & Estado completo del tablero & Cadena de 9 celdas con s\'imbolos \_ / X / O. \\
FTJ / FTU\textbackslash n & Evento de fin de turno & Cierre de turno del robot / del rival. \\
FIN res\textbackslash n & VIC / DER / EMP & Cierre de partida emitido por el robot y consumido por Simulink para estado final. \\
\bottomrule
\end{tabular}
\caption{Formato general de los mensajes intercambiados entre PC central y robot.}
\end{table}

\begin{table}[H]
\centering
\footnotesize
\begin{tabular}{P{0.22\textwidth} P{0.26\textwidth} P{0.42\textwidth}}
\toprule
\textbf{Mensaje} & \textbf{Origen} & \textbf{Momento de uso} \\
\midrule
INI modo & PC central & Justo antes de comenzar la partida. \\
DEC & Robot & Cuando la IA ha calculado su movimiento. \\
ACT cadena9 & Robot & Tras capturar la imagen al terminar el turno rival. \\
FTJ & Robot & Al finalizar el movimiento físico propio. \\
FTU & Robot & Al cerrar el turno del rival humano. \\
FIN res & Robot & Emitido al detectarse victoria, derrota o empate para cerrar la partida. \\
\bottomrule
\end{tabular}
\caption{Secuencia funcional de mensajes empleada durante una partida completa.}
\end{table}

\begin{table}[H]
\centering
\footnotesize
\begin{tabular}{P{0.15\textwidth} P{0.21\textwidth} P{0.56\textwidth}}
\toprule
\textbf{Direcci\'on} & \textbf{Mensaje} & \textbf{Informaci\'on funcional transportada} \\
\midrule
PC central $\rightarrow$ robot & INI modo & Inicio de partida y selecci\'on de modo de operaci\'on (MAN/AUT). \\
PC central $\rightarrow$ robot & MOV origen destino & Orden de recogida y colocaci\'on de ficha en modo manual. \\
Robot $\rightarrow$ PC central & DEC & Confirmaci\'on de decisi\'on tomada en modo autom\'atico. \\
Robot $\rightarrow$ PC central & ACT cadena9 & Estado completo del tablero con codificaci\'on \_ / X / O. \\
Robot $\rightarrow$ PC central & FTJ / FTU & Fin de turno del robot / fin de turno del rival. \\
Robot $\rightarrow$ PC central & FIN res & Resultado final de partida (VIC, DER o EMP), emitido en el flujo operativo actual. \\
\bottomrule
\end{tabular}
\caption{Tabla de interfaz del protocolo de mensajer\'ia PC central--robot.}
\end{table}

\subsubsection{Mensajer\'ia de audio}
La mensajer\'ia de audio forma parte del flujo activo de \texttt{ComRobot.py} mediante llamadas \texttt{robot.say} en hitos clave (inicio de partida, cesi\'on de turno y resultado final), lo que aporta realimentaci\'on local al operador sin depender exclusivamente de la interfaz del PC central.

\subsubsection{IA y decisi\'on autom\'atica}
La toma de decisi\'on en modo autom\'atico implementada en \texttt{ComRobot.py} se concentra en la funci\'on \texttt{decidir\_jugada(fichas\_propias, fichas\_rival)}. El algoritmo empieza calculando las casillas libres a partir de la diferencia entre las fichas propias, las del rival y el tablero completo. Con ese estado construye una lista de candidatos distinta seg\'un la fase de juego: si el robot todav\'ia tiene menos de tres fichas en tablero, los movimientos se generan desde el almac\'en correspondiente (\texttt{alm1}, \texttt{alm2} o \texttt{alm3}) hacia cada casilla libre; cuando ya tiene tres fichas colocadas, los candidatos se obtienen moviendo cada ficha propia hacia cualquiera de las casillas libres.

Antes de elegir al azar, el c\'odigo recorre todos esos candidatos y simula el resultado sobre una copia de las fichas propias para comprobar si la jugada produce victoria inmediata mediante \texttt{verificar\_victoria}. Si existe al menos un movimiento ganador, ese movimiento se devuelve de forma directa. Si no hay victoria inmediata, el algoritmo aplica una regla adicional: en la primera jugada propia, si la casilla 5 est\'a libre, se devuelve la jugada \texttt{alm1 5}; en el resto de situaciones, se selecciona una pareja \texttt{origen--destino} al azar entre todos los movimientos generados. Si no hay candidatos disponibles, la funci\'on devuelve \texttt{(None, None)} y el bucle principal interpreta ese caso como empate. El resultado es una IA muy simple, de coste lineal, pero con una preferencia clara por cerrar partida cuando ya existe una victoria posible.
\section{Validaci\'on de requisitos}

\subsection{Requisitos verificados previamente}
La tabla siguiente refleja el estado de los requisitos ya verificados en anteriores documentos técnicos.

{\footnotesize
\begin{longtable}{@{}>{\raggedright\arraybackslash}p{0.12\textwidth} >{\raggedright\arraybackslash}p{0.15\textwidth} >{\raggedright\arraybackslash}p{0.17\textwidth} >{\raggedright\arraybackslash}p{0.46\textwidth}@{}}
\toprule
\textbf{RqId} & \textbf{Estado actual} & \textbf{Fuente de evidencia} & \textbf{Evidencia t\'ecnica observada} \\
\midrule
\endhead
000C-003 & Verificado & Subsistema de visi\'on del robot & Captura de imagen comprimida, extracci\'on workspace y segmentaci\'on HSV. \\
000C-025 & Verificado & ComRobot.py + SimulinkReceptor.m & Comunicaci\'on remota TCP/IP y parseo de eventos discretos. \\
000C-026 & Verificado & ComRobot.py & Intercambio PC--robot con sockets, mensajes TX/RX y terminador \texttt{\textbackslash n}. \\
000C-030 & Verificado & Inicializaci\'on del control del robot & L\'imite de velocidad de brazo fijado a 50 en la configuraci\'on de arranque. \\
000C-034 & Verificado & ComRobot.py & Fin de turno rival por lectura de entrada digital DI1. \\
000C-038 & Verificado & ComRobot.py & Mensajes codificados como ASCII/UTF-8 con salto de l\'inea final. \\
000C-009 & Verificado & ComRobot.py & Modo manual activo con recepci\'on MOV, sin validaciones de robustez completas. \\
000C-010 & Verificado & ComRobot.py & Modo autom\'atico activo con heur\'istica aleatoria; no minimax en versi\'on actual. \\
000C-036 & Verificado & ComRobot.py & La funci\'on de victoria existe, pero no cierra el bucle principal de partida. \\
000C-037 & Verificado & ComRobot.py & MAX\_TURNOS definido, pendiente de aplicaci\'on efectiva en flujo principal. \\
\bottomrule
\end{longtable}     
}

\subsection{Verificación de requisitos del subsistema}

Se han realizado varias partidas completas para los modos de juego de las cuales se han elegido dos videos. Los enlaces dirigen al momento del video donde se muestra la verificación concreta de cada requisito.

\subsubsection{Sistema de juego 3 en raya}
\textbf{RqId:} [000C-001]
\vspace{0.3em}
\par [Descripción del sistema de juego 3 en raya]
\\ \textbf{Verificación:} Test - Ejecutar partidas completas y verificar cumplimiento de reglas en todos los modos habilitados.
\\ \textbf{Resultado:} Se completan dos partidas: una en modo manual (\url{https://youtu.be/pjByXUHJplQ}) y otra en modo automático (\url{https://youtu.be/cVZLd5qQ-9E}).


\subsubsection{Movimiento por turno}
\textbf{RqId:} [000C-005]
\vspace{0.3em}
\par Verificar que el robot solo se mueve en su turno y permanece en espera en turno rival.
\\ \textbf{Verificación:} Test - Observar comportamiento del robot durante partidas de prueba alternando turnos.
\\ \textbf{Resultado:} Hasta que no se toma la decisión en modo manual el robot no se mueve (\url{https://www.youtube.com/watch?v=pjByXUHJplQ&t=11s}) y hasta que no se pulsa el botón, no inicia el robot el movimiento en modo automático (\url{https://youtu.be/cVZLd5qQ-9E&t=48s}).
\subsubsection{Piezas apiladas - validación en tablero}
\textbf{RqId:} [000C-007]
\vspace{0.3em}
\par Validar uso de pila mientras existan piezas y reubicación desde tablero al agotarse.
\\ \textbf{Verificación:} Test - Ejecutar partida completa y registrar acceso a almacenes y traslados.
\\ \textbf{Resultado:} En la partida en modo automático se prioriza siempre la pieza superior del almacén (\url{https://youtu.be/cVZLd5qQ-9E}). El código que lo demuestra es el siguiente:
\begin{lstlisting}[style=pythoncode, caption={Generación de movimientos priorizando piezas del almacén.}]
if len(fichas_propias) < 3:
    origen_alm = f"alm{len(fichas_propias) + 1}"
    for libre in libres:
        movimientos.append((origen_alm, str(libre)))
else:
    for origen in fichas_propias:
        for libre in libres:
            movimientos.append((str(origen), str(libre)))
\end{lstlisting}

\subsubsection{Modos de juego}
\textbf{RqId:} [000C-008]
\vspace{0.3em}
\par Comprobar disponibilidad y seleción de modos (manual, automático, competición) en interfaz y máquina de estados.
\\ \textbf{Verificación:} Test - Verificar que los tres modos son seleccionables y funcionan de forma independiente.
\\ \textbf{Resultado:} Los modos manual y automático están disponibles en el interfaz de Simulink (\url{https://youtu.be/cVZLd5qQ-9E&t=3s}). El modo competición está de manera experimental pero no implementado en el interfaz ya que se podrá jugar en manual o automático contra otro robot.

\subsubsection{Modo manual}
\textbf{RqId:} [000C-009]
\vspace{0.3em}
\par Ejecutar batería de partidas completas en modo manual para cerrar validación formal de robustez y errores de entrada.
\\ \textbf{Verificación:} Test de sistema - Realizar al menos 5 partidas completas en modo manual.
\\ \textbf{Resultado:} Se han realizado 5 partidas en modo manual como por ejemplo la siguiente (\url{https://youtu.be/pjByXUHJplQ}).

\subsubsection{Modo automático}
\textbf{RqId:} [000C-010]
\vspace{0.3em}
\par Registrar decisiones automáticas y evaluar calidad de jugada de la heurística actual en campo.
\\ \textbf{Verificación:} Análisis + Test - Ejecutar partidas en modo automático y registrar decisiones.
\\ \textbf{Resultado:} Se han realizado 5 partidas en modo automático como por ejemplo la siguiente (\url{https://youtu.be/cVZLd5qQ-9E}).


\subsubsection{Modo competición}
\textbf{RqId:} [000C-011]
\vspace{0.3em}
\par Ejecutar partida robot vs robot sin intervención manual.
\\ \textbf{Verificación:} Test - Verificar funcionamiento del modo competición con dos robots.
\\ \textbf{Resultado:} Queda pendiente de verificar de manera presencial.

\subsubsection{Posición de inicio}
\textbf{RqId:} [000C-012]
\vspace{0.3em}
\par Confirmar permanencia del robot en home hasta que se pulse Start.
\\ \textbf{Verificación:} Test - Verificar que el robot permanece en posición home antes de iniciar partida.
\\ \textbf{Resultado:} El robot se mantiene en posición de home hasta que no se pulsa el botón de la interfaz que inicia la partida (\url{https://youtu.be/cVZLd5qQ-9E&t=5s}).

\subsubsection{Selección jugador inicial}
\textbf{RqId:} [000C-014]
\vspace{0.3em}
\par Verificar selector y correspondencia con primer turno de la partida.
\\ \textbf{Verificación:} Test - Ejecutar partidas iniciadas por ambos jugadores y verificar que comienza quien corresponde.
\\ \textbf{Resultado:} No se ha implementado dado que siempre empieza el jugador/ordenador.

\subsubsection{Cambio de modo restringido}
\textbf{RqId:} [000C-015]
\vspace{0.3em}
\par Probar cambio de modo en partida y fuera de partida; validar bloqueo/permitido según estado.
\\ \textbf{Verificación:} Test - Intentar cambiar modo durante partida (debe bloquearse) y fuera de partida (debe permitirse).
\\ \textbf{Resultado:} Una vez seleccionado el modo de juego no es posible cambiarlo (\url{https://youtu.be/cVZLd5qQ-9E&t=3s}).

\subsubsection{Desacoplo cinemático Z vs X-Y}
\textbf{RqId:} [000C-017]
\vspace{0.3em}
\par Revisar trayectorias del robot durante recogida y colocación de piezas para verificar desacoplo cinemático.
\\ \textbf{Verificación:} Test - Observar ejecución real: aproximación horizontal, descenso vertical, recogida, ascenso vertical, desplazamiento horizontal.
\\ \textbf{Resultado:} Se establecen dos planos XY de movimiento horizontal entre los que solo hay desplazamientos puramente verticales (\url{https://youtu.be/cVZLd5qQ-9E&t=48s}).

\subsubsection{Recogida de piezas}
\textbf{RqId:} [000C-018]
\vspace{0.3em}
\par Comprobar aproximación X-Y y descenso vertical para agarre de piezas.
\\ \textbf{Verificación:} Test - Verificar secuencia de aproximación horizontal y descenso vertical durante recogida.
\\ \textbf{Resultado:} Se establece una posición de aproximación al almacen de fichas desde la que se accede a las mismas de manera estrictamente vertical (\url{https://youtu.be/cVZLd5qQ-9E&t=80s})

\subsubsection{Ascenso tras recogida}
\textbf{RqId:} [000C-019]
\vspace{0.3em}
\par Verificar ascenso vertical tras agarre antes de cualquier desplazamiento horizontal.
\\ \textbf{Verificación:} Test - Observar que el robot asciende verticalmente tras recoger antes de moverse horizontalmente.
\\ \textbf{Resultado:} Se establecen posiciones de aproximación para todas las posiciones del tablero para que se ascienda verticalmente después del agarre de una ficha (\url{https://youtu.be/pjByXUHJplQ&t=173s}).

\subsubsection{Movimientos X-Y en plano de aproximación}
\textbf{RqId:} [000C-020]
\vspace{0.3em}
\par Confirmar que los desplazamientos horizontales solo se realizan en el plano de aproximación.
\\ \textbf{Verificación:} Test - Verificar que los movimientos horizontales se ejecutan siempre en la altura de aproximación.
\\ \textbf{Resultado:} Todos los puntos que se encuentran en el plano de aproximación están a la misma altura (\url{https://youtu.be/pjByXUHJplQ&t=177s})

\subsubsection{Mensajes de audio}
\textbf{RqId:} [000C-022]
\vspace{0.3em}
\par Integrar audio en flujo principal y validar reprodución en estados clave.
\\ \textbf{Verificación:} Test - Ejecutar partida y verificar mensajes de audio en: inicio, movimiento y fin de partida.
\\ \textbf{Resultado:} Se han establecido mensajes de inicio de partida, turno del rival y final de partida (\url{https://youtu.be/cVZLd5qQ-9E&t=5s}).

\subsubsection{Actualización por visión}
\textbf{RqId:} [000C-023]
\vspace{0.3em}
\par Contrastar tras cada turno el estado detectado con el tablero real.
\\ \textbf{Verificación:} Test - Ejecutar partidas y verificar que la detección visual coincide con el estado real del tablero.
\\ \textbf{Resultado:} Se reciben mensajes de actualización de pantalla después de cada turno jugado que se ven reflejado en el interfaz de Simulink (\url{https://youtu.be/pjByXUHJplQ&t=150s}). Nota: el tablero se representa desde el punto de vista del robot.

\subsubsection{Interfaz de usuario PC central}
\textbf{RqId:} [000C-024]
\vspace{0.3em}
\par Verificar presencia y funcionamiento de elementos requeridos de interfaz en el PC central.
\\ \textbf{Verificación:} Test - Verificar elementos de interfaz: selector de modo, botón Start, visualización de tablero, estado de partida.
\\ \textbf{Resultado:} La interfaz incluye un botón de inicio de partida, un selector de modo de juego, un display del estado de la partida, el tablero de juego actualizado, un display para determinar el turno en modo manual y selectores de posición de origen y destino junto con botón de toma de decisión (\url{https://youtu.be/pjByXUHJplQ&t=195s}).

\subsubsection{Estado del robot conocido}
\textbf{RqId:} [000C-027]
\vspace{0.3em}
\par Revisar trazas para comprobar disponibilidad continua del estado del robot.
\\ \textbf{Verificación:} Análisis - Registrar trazas de comunicación y verificar que el estado es conocido en todo momento.
\\ \textbf{Resultado:} Se dispone de un display que indica si el robot esta en movimiento o en el turno del rival ('Espera') o esperando a la decisión del jugador ('Te toca jugar') (\url{https://youtu.be/pjByXUHJplQ&t=195s}).

\subsubsection{Formato de mensajes PC $\rightarrow$ robot}
\textbf{RqId:} [000C-028]
\vspace{0.3em}
\par Validar estructura y argumentos de mensajes de comando en trazas reales.
\\ \textbf{Verificación:} Test - Capturar trazas de mensajes INI y MOV, verificar que cumplen formato especificado.
\\ \textbf{Resultado:} [Completar tras verificación]

\subsubsection{Formato de mensajes robot $\rightarrow$ PC}
\textbf{RqId:} [000C-029]
\vspace{0.3em}
\par Validar estructura y argumentos de mensajes de estado/evento en trazas reales.
\\ \textbf{Verificación:} Test - Capturar trazas de mensajes DEC, ACT, FTJ, FTU, FIN; verificar que cumplen formato especificado.
\\ \textbf{Resultado:} [Completar tras verificación]

\subsubsection{Tiempo máximo de recogida}
\textbf{RqId:} [000C-031]
\vspace{0.3em}
\par Medir al menos 5 ejecuciones de recogida y verificar que el tiempo es $\leq 15$ segundos.
\\ \textbf{Verificación:} Test - Ejecutar 5 recogidas consecutivas, cronometrar y registrar tiempos.
\\ \textbf{Resultado:} Se realizan varias ciclos de medición de tiempos de recogida de piezas, en el siguiente video se puede ver que desde que recibe la orden ('1:50') hasta que agarra la pieza ('1:59') pasan 9 segundos (\url{https://youtu.be/pjByXUHJplQ&t=110s}).

\subsubsection{Tiempo máximo de transporte}
\textbf{RqId:} [000C-032]
\vspace{0.3em}
\par Medir al menos 5 transportes completos y verificar que el tiempo es $\leq 15$ segundos.
\\ \textbf{Verificación:} Test - Ejecutar 5 ciclos de transporte completo, cronometrar y registrar tiempos.
\\ \textbf{Resultado:} Se realizan varias ciclos de medición de tiempos desde que se agarra la pieza hasta que se suelta. en el siguiente video se puede ver que pasan 10 segundos desde que se agarra ('1:59') hasta que se suelta ('2:09') (\url{https://youtu.be/pjByXUHJplQ&t=119s}).

\subsubsection{Tiempo máximo de decisión automática}
\textbf{RqId:} [000C-033]
\vspace{0.3em}
\par Medir tiempo de decisión en modo automático/competición y verificar que es $\leq 10$ segundos.
\\ \textbf{Verificación:} Test - Ejecutar partidas en modo automático, medir tiempos de decisión y registrar.
\\ \textbf{Resultado:} La decisión en modo automático se toma de manera instantánea ya que cuando el rival pulsa el botón físico de final de turno el robot inicia su movimiento (\url{https://youtu.be/cVZLd5qQ-9E&t=48s}).

\subsubsection{Fin de turno rival por botón físico}
\textbf{RqId:} [000C-034]
\vspace{0.3em}
\par Accionar botón físico (entrada digital DI1) y comprobar transición al siguiente estado.
\\ \textbf{Verificación:} Test - Ejecutar batería repetida de accionamientos de botón físico, verificar que se detecta correctamente.
\\ \textbf{Resultado:} El rival determina el final de su turno mediante la pulsación del botón físico conectado a la entrada digital 1 del robot (\url{https://youtu.be/cVZLd5qQ-9E&t=48s}).

\section{Gu\'ia r\'apida para la puesta en marcha}

\subsection{Prerrequisitos}
\textbf{Hardware m\'inimo:}
\begin{itemize}
  \item Robot Ned-2 calibrado.
  \item C\'amara del robot operativa.
  \item PC central y PC de control del robot conectados por red.
  \item Tablero y piezas en posici\'on inicial.
\end{itemize}

\textbf{Software m\'inimo:}
\begin{itemize}
  \item Entorno Python con dependencias del proyecto.
  \item Modelo Simulink con bloques de emisi\'on/recepci\'on configurados.
  \item Scripts de control de robot y visi\'on disponibles.
\end{itemize}

\subsection{Secuencia de arranque recomendada}
\begin{enumerate}
  \item Verificar conexi\'on de red entre PC central y PC del robot.
  \item Encender y calibrar el Ned-2.
  \item Cargar configuraci\'on de poses y par\'ametros de visi\'on.
  \item Iniciar servicios/scripts de control en el PC del robot.
  \item Abrir el modelo de control en PC central (Simulink).
  \item Configurar IP/puertos y arrancar la comunicaci\'on.
  \item Confirmar estado home y seleccionar modo de juego.
  \item Pulsar Start para iniciar partida.
\end{enumerate}

\section{Conclusiones}

\end{document}